\documentclass[12pt,letterpaper]{report}
\usepackage{graphicx}
\usepackage{scrextend}
\usepackage{vmargin}
\usepackage{graphicx}
\usepackage{multirow}
\usepackage[utf8]{inputenc}
\usepackage[spanish]{babel}
\usepackage{multicol}
\usepackage{enumerate}
\usepackage{float}
\usepackage{amsmath, amsthm, amssymb, amsfonts}
\usepackage[usenames]{color}
\parindent=0mm
\pagestyle{empty}
\definecolor{miorange}{rgb}{0.91, 0.43, 0.0}
\begin{document}
\setmargins{2.5cm}      
{1.5cm}                     
{2cm}  
{24cm}                    
{10pt}                          
{1cm}                          
{0pt}                             
{2cm}
\begin{flushright}
\textbf{\today} $TE$
\end{flushright}
\vspace{-1.8cm}
\section*{Organización de las clases}
\subsection*{Clase 1}
\subsubsection*{Objetivo:}
Reconocer de manera clara las propiedades que se encuentran en el grupo de los números reales, aplicar operaciones de suma, resta, multiplicación, división y potencia de números reales ejecutando el orden de prioridad que tiene cada operación aritmética .
\subsubsection*{Temas:}
\begin{enumerate}
    \item Propiedades de los números reales.
    \item Operaciones con números reales. 
    \item Simplificación de números.
    \item Jerarquía de operaciones.
\end{enumerate}
\subsection*{Clase 2}
\subsubsection*{Objetivo:}
Reconocer y aplicar de manera precisa operaciones algebraicas.
\subsubsection*{Temas:}
\begin{enumerate}
    \item Suma y resta de polinomios.
    \item Multiplicación de monomios.
    \item División de monomios.
    \item Potencia de monomios.
\end{enumerate}
\subsection*{Clase 3}
\subsubsection*{Objetivo:}
Realizar de manera eficiente la operación de multiplicación con polinomios.
\subsubsection*{Temas:}
\begin{enumerate}
    \item Multiplicación de polinomios.
\end{enumerate}
\pagebreak
\subsection*{Clase 4}
\subsubsection*{Objetivo:}
Reconocer y aplicar una factorización a una ecuación cuadrática o términos comunes.
\subsubsection*{Temas:}
\begin{enumerate}
    \item Factorización de ecuaciones cuadráticas de la forma $ax^2+bx+c$.
    \item Factorización de términos comunes.
\end{enumerate}
\subsection*{Clase 5}
\subsubsection*{Objetivo:}
Realizar de manera eficiente la simplificación de expresiones algebraicas largas.
\subsubsection*{Temas:}
\begin{enumerate}
    \item Suma y resta de fracciones algebraicas de polinomios.
    \item Multiplicación y división de fracciones algebraicas de polinomios.
\end{enumerate}
\subsection*{Clase 6}
\subsubsection*{Objetivo:}
Reconocer y aplicar las diferentes maneras en las cuales puedes obtener la solución a una ecuación.
\subsubsection*{Temas:}
\begin{enumerate}
    \item Solución a funciones y/o expresiones algebraicas.
\end{enumerate}
\subsection*{Clase 7}
\subsubsection*{Objetivo:}
Resolver sistemas de ecuaciones lineales dadas las ecuaciones.
\subsubsection*{Temas:}
\begin{enumerate}
    \item Método suma y resta/sustitución/igualación.
\end{enumerate}
\textbf{Nota:} Se tomara el método que se les sea más sencillo realizar y analizar, ya que lo que se busca es que realicen la solución a partir de su lógica. 
\pagebreak
\subsection*{Clase 8}
\subsubsection*{Objetivo:}
Plantear y resolver un problema el cual involucre resolver un sistema de ecuaciones.
\subsubsection*{Temas:}
\begin{enumerate}
    \item Sistemas de ecuaciones razonados.
\end{enumerate}
\subsection*{Clase 9}
\subsubsection*{Objetivo:}
Resolver problemas razonados aplicando álgebra.
\subsubsection*{Temas:}
\begin{enumerate}
    \item Aplicación de factorización.
    \item Planteamiento de expresiones algebraicas en áreas y volúmenes.
\end{enumerate}
\subsection*{Clase 10}
\subsubsection*{Objetivo:}
Comprender y ubicar los conceptos de los triángulos y aplicar proporciones dentro de el.
\subsubsection*{Temas:}
\begin{enumerate}
    \item Propiedades de la semejanza en triángulos y sus proporciones.
\end{enumerate}
\subsection*{Clase 11}
\subsubsection*{Objetivo:}
Aplicar de manera efectiva el teorema de Pitágoras.
\subsubsection*{Temas:}
\begin{enumerate}
    \item Teorema de Pitágoras.
\end{enumerate}
\subsection*{Clase 12}
\subsubsection*{Objetivo:}
Reconocer y aplicar las funciones periódicas-
\subsubsection*{Temas:}
\begin{enumerate}
    \item Función $\sin{x}, \cos{x}, \tan{x}$.
\end{enumerate}
\subsection*{Clase 13}
\subsubsection*{Objetivo:}
Aplicar de manera efectiva los teoremas dentro de un triangulo para obtener toda la información de el.
\subsubsection*{Temas:}
\begin{enumerate}
\item Resolver triángulos por medio de funciones periódicas y teorema de Pitágoras.
\end{enumerate}
\subsection*{Clase 14}
\subsubsection*{Objetivo:}
Reconocer la diferencia entre permutaciones y combinaciones, calcular ya sea el caso.
\subsubsection*{Temas:}
\begin{enumerate}
    \item Permutaciones
    \item Combinaciones
\end{enumerate}
\subsection*{Clase 15}
\subsubsection*{Objetivo:}
Calcular la probabilidad de sucesos independientes, ya sea varios sucesos consecutivos o sucesos simultáneos.
\subsubsection*{Temas:}
\begin{enumerate}
    \item Probabilidad de sucesos independientes.
    \item Probabilidad de sucesos independientes consecutivos o simultáneos.
\end{enumerate}
\end{document}